In Appendice vengono presentati ulteriori esempi di utilizzo delle tecniche illustrate nei Capitoli precedenti, nonché le soluzioni ad alcuni dei più famosi problemi di Sincronizzazione tra processi.

\section{Problema dei Processi Lettori e Scrittori} \label{appendix:ReadersAndWriters}
Si riporta a seguire un'implementazione della soluzione al problema dei Processi Lettori e Scrittori tramite l'utilizzo dei \texttt{ReentrantReadWriteLock}.
\input{images/ReadersAndWriters.tex}

\section{Problema della Cena tra i Cinque Filosofi} \label{appendix:FivePhilosophers}
Si riporta a seguire un'implementazione della soluzione al problema della Cena tra i Cinque Filosofi tramite l'utilizzo dei Semafori.
\begin{lstlisting}[language=Java, numbers=left, stepnumber=1]
/*
 * File: DiningPhilosophers.java
 * Il Programma mira ad implementare una soluzione del classico problema
 * della cena tra i Filosofi tramite l'utilizzo delle metodologie
 * di sincronizzazione messe a disposizione dal Linguaggio Java
 */

import java.util.Random;
import java.util.concurrent.Semaphore;

public class DiningPhilosophers {
    
    // Array di semafori per l'acquisizione delle bacchette
    private Semaphore[] sticksSemaphores = new Semaphore[5];

    public static void main(String[] args) {
        new DiningPhilosophers();
    }

    private DiningPhilosophers() {
        // Inizializzo i Semafori come mutualmente esclusivi ed equi
        for(int i=0; i<5; i++) {
            sticksSemaphores[i] = new Semaphore(1, true);
        }

        // Creo ed avvio i Filosofi
        for(int i=0; i<5; i++) {
            DiningPhilosophersPhilosopher philosopherRunnable =
                new DiningPhilosophersPhilosopher(i+1, sticksSemaphores);
            Thread philosopherThread = new Thread(philosopherRunnable);
            philosopherThread.start();
        }
    }
}

class DiningPhilosophersPhilosopher implements Runnable {

    // Identificativo del filosofo
    private int philosopherNum;
    // Semafori per l'acquisizione delle bacchette
    private Semaphore[] stickSemaphores;

    // Costruttore
    protected DiningPhilosophersPhilosopher(int philosopherNum, Semaphore[] sticksSemaphores) {
        this.philosopherNum = philosopherNum;
        this.stickSemaphores = sticksSemaphores;
        System.out.println("Filosofo numero " + this.philosopherNum + " creato");
    }

    // Metodo per far pensare il filosofo 
    // aggiungo come parametro se sta pensando perche' non ha mangiato per
    // farlo mangiare un numero fissato di volte
    private void think(boolean becauseDidntEat) {
        // Scelgo un tempo per pensare compreso tra 200 e 1000ms
        int thinkTime = new Random(System.currentTimeMillis()).nextInt(200, 1001);
        System.out.println("Il filosofo numero " + this.philosopherNum +
            " pensera' per " + thinkTime/1000.0 + " secondi");
        try {
            // Simulo il pensiero
            Thread.sleep(thinkTime);
        } catch(InterruptedException ie) {
            System.out.println("Il Thread e' stato interrotto inaspettatamente");
            ie.printStackTrace();
        }
        System.out.println("Il filosofo numero " + this.philosopherNum +
            " ha terminato di pensare");
        // Se sono entrato in think() perche' non ho mangiato, torno su eat()
        if(becauseDidntEat) {
            eat();
        }
    }

    // Metodo per far mangiare il filosofo
    // Il controllo della disponibilita' delle bacchette va effettuato in una sezione critica
    private synchronized void eat() {
        System.out.println("Il filosofo numero " + this.philosopherNum +
            " tenta di acquisire le bacchette");
        // Controllo se entrambe le bacchette sono disponibili
        if(this.stickSemaphores[(philosopherNum-1)%5].availablePermits() == 1
            && this.stickSemaphores[(philosopherNum)%5].availablePermits() == 1) {
            try {
                // Acquisizione Bacchetta sinistra
                this.stickSemaphores[(philosopherNum-1)%5].acquire();
                // Acquisizione Bacchetta destra
                this.stickSemaphores[(philosopherNum)%5].acquire();
                // Mangia
                System.out.println("Il filosofo " + this.philosopherNum + " sta mangiando");
                Thread.sleep(100);
            } catch(InterruptedException ie) {
                System.out.println("Il Thread e' stato interrotto inaspettatamente");
                ie.printStackTrace();
            } finally {
                // Rilascio Bacchetta sinistra
                this.stickSemaphores[(philosopherNum-1)%5].release();
                // Rilascio Bacchetta destra
                this.stickSemaphores[(philosopherNum)%5].release();
                // Termina di mangiare
                System.out.println("Il filosofo " + this.philosopherNum + " ha terminato di mangiare");
            }
        } else {
            System.out.println("Il filosofo numero " + this.philosopherNum +
            " non e' riuscito nell'acquisizione delle bacchette, pensera' e tornera' per mangiare");
            // Chiamo think() con l'argomento true
            think(true);
        }
    }

    // Implementazione di run
    @Override
    public void run() {
        System.out.println("Filosofo numero " + this.philosopherNum + " avviato");
        for(int i=0; i<3; i++) {
            eat();
            think(false);
        }
    }
}
\end{lstlisting}
\captionof{lstlisting}{Programma Problema della Cena tra i 5 Filosofi}
\begin{lstlisting}[language=Bash]
$ java DiningPhilosophers
Filosofo numero 1 creato
Filosofo numero 2 creato
Filosofo numero 3 creato
Filosofo numero 1 avviato
Filosofo numero 3 avviato
Filosofo numero 4 creato
Filosofo numero 2 avviato
Il filosofo numero 1 tenta di acquisire le bacchette
Filosofo numero 5 creato
Il filosofo numero 3 tenta di acquisire le bacchette
Il filosofo numero 2 tenta di acquisire le bacchette
Filosofo numero 4 avviato
Filosofo numero 5 avviato
Il filosofo 3 sta mangiando
Il filosofo 1 sta mangiando
Il filosofo numero 5 tenta di acquisire le bacchette
Il filosofo numero 5 non e' riuscito nell'acquisizione delle bacchette, pensera' e tornera' per mangiare
Il filosofo numero 4 tenta di acquisire le bacchette
Il filosofo numero 2 non e' riuscito nell'acquisizione delle bacchette, pensera' e tornera' per mangiare
Il filosofo numero 4 non e' riuscito nell'acquisizione delle bacchette, pensera' e tornera' per mangiare
Il filosofo numero 2 pensera' per 0.962 secondi
Il filosofo numero 4 pensera' per 0.962 secondi
Il filosofo numero 5 pensera' per 0.962 secondi
Il filosofo 3 ha terminato di mangiare
Il filosofo 1 ha terminato di mangiare
Il filosofo numero 3 pensera' per 0.31 secondi
Il filosofo numero 1 pensera' per 0.438 secondi
Il filosofo numero 3 ha terminato di pensare
Il filosofo numero 3 tenta di acquisire le bacchette
Il filosofo 3 sta mangiando
Il filosofo 3 ha terminato di mangiare
Il filosofo numero 3 pensera' per 0.795 secondi
Il filosofo numero 1 ha terminato di pensare
Il filosofo numero 1 tenta di acquisire le bacchette
Il filosofo 1 sta mangiando
Il filosofo 1 ha terminato di mangiare
Il filosofo numero 1 pensera' per 0.382 secondi
Il filosofo numero 2 ha terminato di pensare
Il filosofo numero 2 tenta di acquisire le bacchette
Il filosofo 2 sta mangiando
Il filosofo numero 4 ha terminato di pensare
Il filosofo numero 5 ha terminato di pensare
Il filosofo numero 4 tenta di acquisire le bacchette
Il filosofo numero 5 tenta di acquisire le bacchette
Il filosofo numero 5 non e' riuscito nell'acquisizione delle bacchette, pensera' e tornera' per mangiare
Il filosofo 4 sta mangiando
Il filosofo numero 5 pensera' per 0.25 secondi
Il filosofo numero 1 ha terminato di pensare
Il filosofo numero 1 tenta di acquisire le bacchette
Il filosofo numero 1 non e' riuscito nell'acquisizione delle bacchette, pensera' e tornera' per mangiare
Il filosofo numero 1 pensera' per 0.977 secondi
Il filosofo 2 ha terminato di mangiare
Il filosofo 4 ha terminato di mangiare
Il filosofo numero 2 pensera' per 0.398 secondi
Il filosofo numero 4 pensera' per 0.264 secondi
Il filosofo numero 5 ha terminato di pensare
Il filosofo numero 5 tenta di acquisire le bacchette
Il filosofo 5 sta mangiando
Il filosofo numero 3 ha terminato di pensare
Il filosofo numero 3 tenta di acquisire le bacchette
Il filosofo 3 sta mangiando
Il filosofo 5 ha terminato di mangiare
Il filosofo numero 5 pensera' per 0.878 secondi
Il filosofo numero 4 ha terminato di pensare
Il filosofo numero 4 tenta di acquisire le bacchette
Il filosofo numero 4 non e' riuscito nell'acquisizione delle bacchette, pensera' e tornera' per mangiare
Il filosofo numero 4 pensera' per 0.564 secondi
Il filosofo 3 ha terminato di mangiare
Il filosofo numero 3 pensera' per 0.63 secondi
Il filosofo numero 2 ha terminato di pensare
Il filosofo numero 2 tenta di acquisire le bacchette
Il filosofo 2 sta mangiando
Il filosofo 2 ha terminato di mangiare
Il filosofo numero 2 pensera' per 0.961 secondi
Il filosofo numero 4 ha terminato di pensare
Il filosofo numero 4 tenta di acquisire le bacchette
Il filosofo 4 sta mangiando
Il filosofo numero 1 ha terminato di pensare
Il filosofo numero 1 tenta di acquisire le bacchette
Il filosofo 1 sta mangiando
Il filosofo 4 ha terminato di mangiare
Il filosofo numero 4 pensera' per 0.877 secondi
Il filosofo numero 3 ha terminato di pensare
Il filosofo 1 ha terminato di mangiare
Il filosofo numero 1 pensera' per 0.622 secondi
Il filosofo numero 5 ha terminato di pensare
Il filosofo numero 5 tenta di acquisire le bacchette
Il filosofo 5 sta mangiando
Il filosofo 5 ha terminato di mangiare
Il filosofo numero 5 pensera' per 0.247 secondi
Il filosofo numero 2 ha terminato di pensare
Il filosofo numero 2 tenta di acquisire le bacchette
Il filosofo 2 sta mangiando
Il filosofo numero 5 ha terminato di pensare
Il filosofo numero 5 tenta di acquisire le bacchette
Il filosofo 5 sta mangiando
Il filosofo 2 ha terminato di mangiare
Il filosofo numero 2 pensera' per 0.272 secondi
Il filosofo 5 ha terminato di mangiare
Il filosofo numero 5 pensera' per 0.226 secondi
Il filosofo numero 1 ha terminato di pensare
Il filosofo numero 5 ha terminato di pensare
Il filosofo numero 4 ha terminato di pensare
Il filosofo numero 4 tenta di acquisire le bacchette
Il filosofo 4 sta mangiando
Il filosofo numero 2 ha terminato di pensare
Il filosofo 4 ha terminato di mangiare
Il filosofo numero 4 pensera' per 0.334 secondi
Il filosofo numero 4 ha terminato di pensare
\end{lstlisting}
\captionof{lstlisting}{Output del Programma di cui sopra.}

\section{Esempio di utilizzo delle Primitive di \texttt{synchronized}} \label{appendix:AlternateCounterSynchronized}
Viene implementato un Contatore condiviso tra due \textit{Thread}. Il \textit{Thread} 1 potrà aggiornare il contatore se e solo se questo è dispari, il \textit{Thread} 2 solo se è pari.
\begin{lstlisting}[language=Java, numbers=left, stepnumber=1]
/*
 * File: AlternateCounterSynchronized.java
 * Viene implementato un Contatore condiviso tra due Thread.
 * Il Thread 1 potra' aggiornare il contatore se e solo se questo e' dispari,
 * il Thread 2 solo se e' pari.
 */

public class AlternateCounterSynchronized {

    // Metodo main per l'avvio
    public static void main(String[] args) {
        new AlternateCounterSynchronized();
    }

    // Costruttore dell'oggetto principale
    private AlternateCounterSynchronized() {
        // Creo l'oggetto del counter da condividere
        AlternateCounterSynchronizedCounter acl = new AlternateCounterSynchronizedCounter();
        // Creo e avvio i due Thread
        AlternateCounterSynchronizedRunnable run1 = new AlternateCounterSynchronizedRunnable(1, acl);
        AlternateCounterSynchronizedRunnable run2 = new AlternateCounterSynchronizedRunnable(2, acl);
        Thread t1 = new Thread(run1);
        Thread t2 = new Thread(run2);
        t1.start();
        t2.start();
        // Attendo la loro terminazione e stampo il risultato
        try {
            t1.join();
            t2.join();
        } catch(InterruptedException ie) {
            System.out.println("Errore nell'attesa dei Thread figli");
            ie.printStackTrace();
        }
        System.out.println("Il valore finale del contatore e': "
            + acl.getCounter());
    }
}

class AlternateCounterSynchronizedCounter {

    // Variabile contatore
    private int counter = 0;

    // Funzione per l'incremento del contatore
    protected synchronized void increaseAndPrintCounter(int threadNum) {
        // Finche' la condizione non viene rispettata, non incrementare
        while(threadNum%2 != this.counter%2) {
            try {
                System.out.println("Sono il Thread " + threadNum + " e sono bloccato in attesa della condizione");
                wait();
            } catch(InterruptedException ie) {
                System.out.println("Thread interrotto inaspettatamente");
                ie.printStackTrace();
            }
        }
        // Incrementa
        counter++;
        // Stampo il risultato dell'operazione
        System.out.println("Sono il Thread " + threadNum + " ed ho aggiornato il contatore: " + this.counter);
        // Segnala all'altro processo eventualmente bloccato
        notify();
    }

    // Metodo di supporto
    protected int getCounter() {
        return this.counter;
    }
}

class AlternateCounterSynchronizedRunnable implements Runnable {

    // Identificativo del Thread
    private int threadNum;
    // Risorsa condivisa
    private AlternateCounterSynchronizedCounter acl;

    // Costruttore
    protected AlternateCounterSynchronizedRunnable(int threadNum, AlternateCounterSynchronizedCounter acl) {
        this.threadNum = threadNum;
        this.acl = acl;
        System.out.println("Runnable " + this.threadNum + " creato");
    }
    
    // Implementazione del Metodo run()
    @Override
    public void run() {
        System.out.println("Thread " + this.threadNum + " avviato");
        for(int i=0; i<10; i++) {
            acl.increaseAndPrintCounter(this.threadNum);
        }
    }
}
\end{lstlisting}
\captionof{lstlisting}{Implementazione tramite l'utilizzo di \textit{Lock} e Variabili Condizione.}
\begin{lstlisting}[language=Bash]
$ java AlternateCounterLock
Runnable 1 creato
Runnable 2 creato
Thread 2 avviato
Thread 1 avviato
Sono il Thread 2 ed ho aggiornato il contatore: 1
Sono il Thread 2 e sono bloccato in attesa della condizione
Sono il Thread 1 ed ho aggiornato il contatore: 2
Sono il Thread 1 e sono bloccato in attesa della condizione
Sono il Thread 2 ed ho aggiornato il contatore: 3
Sono il Thread 2 e sono bloccato in attesa della condizione
Sono il Thread 1 ed ho aggiornato il contatore: 4
Sono il Thread 1 e sono bloccato in attesa della condizione
Sono il Thread 2 ed ho aggiornato il contatore: 5
Sono il Thread 2 e sono bloccato in attesa della condizione
Sono il Thread 1 ed ho aggiornato il contatore: 6
Sono il Thread 1 e sono bloccato in attesa della condizione
Sono il Thread 2 ed ho aggiornato il contatore: 7
Sono il Thread 2 e sono bloccato in attesa della condizione
Sono il Thread 1 ed ho aggiornato il contatore: 8
Sono il Thread 1 e sono bloccato in attesa della condizione
Sono il Thread 2 ed ho aggiornato il contatore: 9
Sono il Thread 2 e sono bloccato in attesa della condizione
Sono il Thread 1 ed ho aggiornato il contatore: 10
Sono il Thread 1 e sono bloccato in attesa della condizione
Sono il Thread 2 ed ho aggiornato il contatore: 11
Sono il Thread 2 e sono bloccato in attesa della condizione
Sono il Thread 1 ed ho aggiornato il contatore: 12
Sono il Thread 1 e sono bloccato in attesa della condizione
Sono il Thread 2 ed ho aggiornato il contatore: 13
Sono il Thread 2 e sono bloccato in attesa della condizione
Sono il Thread 1 ed ho aggiornato il contatore: 14
Sono il Thread 1 e sono bloccato in attesa della condizione
Sono il Thread 2 ed ho aggiornato il contatore: 15
Sono il Thread 2 e sono bloccato in attesa della condizione
Sono il Thread 1 ed ho aggiornato il contatore: 16
Sono il Thread 1 e sono bloccato in attesa della condizione
Sono il Thread 2 ed ho aggiornato il contatore: 17
Sono il Thread 2 e sono bloccato in attesa della condizione
Sono il Thread 1 ed ho aggiornato il contatore: 18
Sono il Thread 1 e sono bloccato in attesa della condizione
Sono il Thread 2 ed ho aggiornato il contatore: 19
Sono il Thread 1 ed ho aggiornato il contatore: 20
Il valore finale del contatore e': 20
\end{lstlisting}
\captionof{lstlisting}{Output del Programma di cui sopra.}

\section{Esempio di utilizzo delle Variabili Condizione} \label{appendix:AlternateCounterConditionVars}
Si riporta una reimplementazione dell'esempio precedente utilizzando le Variabili Condizione di un \textit{Lock}.
\begin{lstlisting}[language=Java, numbers=left, stepnumber=1]
/*
 * File: AlternateCounterLock.java
 * Viene implementato un Contatore condiviso tra due Thread.
 * Il Thread 1 potra' aggiornare il contatore se e solo se questo e' dispari,
 * il Thread 2 solo se e' pari.
 */

import java.util.concurrent.locks.Condition;
import java.util.concurrent.locks.Lock;
import java.util.concurrent.locks.ReentrantLock;

public class AlternateCounterLock {

    // Metodo main per l'avvio
    public static void main(String[] args) {
        new AlternateCounterLock();
    }

    // Costruttore dell'oggetto principale
    private AlternateCounterLock() {
        // Creo l'oggetto del counter da condividere
        AlternateCounterLockCounter acl = new AlternateCounterLockCounter();
        // Creo e avvio i due Thread
        AlternateCounterLockRunnable run1 = new AlternateCounterLockRunnable(1, acl);
        AlternateCounterLockRunnable run2 = new AlternateCounterLockRunnable(2, acl);
        Thread t1 = new Thread(run1);
        Thread t2 = new Thread(run2);
        t1.start();
        t2.start();
        // Attendo la loro terminazione e stampo il risultato
        try {
            t1.join();
            t2.join();
        } catch(InterruptedException ie) {
            System.out.println("Errore nell'attesa dei Thread figli");
            ie.printStackTrace();
        }
        System.out.println("Il valore finale del contatore e': "
            + acl.getCounter());
    }
}

class AlternateCounterLockCounter {

    // Variabile contatore
    private int counter = 0;
    // Lock per l'accesso in mutua esclusione alla sezione critica
    // non equo per aumentare la possibilita' di bloccarsi sulla condizione
    private Lock lck = new ReentrantLock(false);
    // Variabile condizione
    private Condition parity = lck.newCondition();

    // Funzione per l'incremento del contatore
    protected void increaseAndPrintCounter(int threadNum) {
        // Acquisizione del Lock
        lck.lock();
        try {
            // Finche' la condizione non viene rispettata, non incrementare
            while(threadNum%2 != this.counter%2) {
                System.out.println("Sono il Thread " + threadNum +
                    " e sono bloccato sulla variabile condizione");
                parity.await();
            }
            // Incrementa
            counter++;
            // Stampo il risultato dell'operazione
            System.out.println("Sono il Thread " + threadNum +
                " ed ho aggiornato il contatore: " + this.counter);
            // Segnala all'altro processo eventualmente bloccato
            parity.signal();
        } catch(InterruptedException ie) {
            System.out.println("Thread interrotto inaspettatamente");
            ie.printStackTrace();
        } finally {
            // Rilascio del Lock
            lck.unlock();
        }
    }

    // Metodo di supporto
    protected int getCounter() {
        return this.counter;
    }
}

class AlternateCounterLockRunnable implements Runnable {

    // Identificativo del Thread
    private int threadNum;
    // Risorsa condivisa
    private AlternateCounterLockCounter acl;

    // Costruttore
    protected AlternateCounterLockRunnable(int threadNum, AlternateCounterLockCounter acl) {
        this.threadNum = threadNum;
        this.acl = acl;
        System.out.println("Runnable " + this.threadNum + " creato");
    }
    
    // Implementazione del Metodo run()
    @Override
    public void run() {
        System.out.println("Thread " + this.threadNum + " avviato");
        for(int i=0; i<10; i++) {
            acl.increaseAndPrintCounter(this.threadNum);
        }
    }
}
\end{lstlisting}
\captionof{lstlisting}{Implementazione tramite l'utilizzo di \textit{Lock} e Variabili Condizione.}
\begin{lstlisting}[language=Bash]
$ java AlternateCounterLock
Runnable 1 creato
Runnable 2 creato
Thread 1 avviato
Thread 2 avviato
Sono il Thread 1 e sono bloccato sulla variabile condizione
Sono il Thread 2 ed ho aggiornato il contatore: 1
Sono il Thread 2 e sono bloccato sulla variabile condizione
Sono il Thread 1 ed ho aggiornato il contatore: 2
Sono il Thread 1 e sono bloccato sulla variabile condizione
Sono il Thread 2 ed ho aggiornato il contatore: 3
Sono il Thread 2 e sono bloccato sulla variabile condizione
Sono il Thread 1 ed ho aggiornato il contatore: 4
Sono il Thread 1 e sono bloccato sulla variabile condizione
Sono il Thread 2 ed ho aggiornato il contatore: 5
Sono il Thread 2 e sono bloccato sulla variabile condizione
Sono il Thread 1 ed ho aggiornato il contatore: 6
Sono il Thread 1 e sono bloccato sulla variabile condizione
Sono il Thread 2 ed ho aggiornato il contatore: 7
Sono il Thread 2 e sono bloccato sulla variabile condizione
Sono il Thread 1 ed ho aggiornato il contatore: 8
Sono il Thread 1 e sono bloccato sulla variabile condizione
Sono il Thread 2 ed ho aggiornato il contatore: 9
Sono il Thread 2 e sono bloccato sulla variabile condizione
Sono il Thread 1 ed ho aggiornato il contatore: 10
Sono il Thread 1 e sono bloccato sulla variabile condizione
Sono il Thread 2 ed ho aggiornato il contatore: 11
Sono il Thread 2 e sono bloccato sulla variabile condizione
Sono il Thread 1 ed ho aggiornato il contatore: 12
Sono il Thread 1 e sono bloccato sulla variabile condizione
Sono il Thread 2 ed ho aggiornato il contatore: 13
Sono il Thread 2 e sono bloccato sulla variabile condizione
Sono il Thread 1 ed ho aggiornato il contatore: 14
Sono il Thread 2 ed ho aggiornato il contatore: 15
Sono il Thread 2 e sono bloccato sulla variabile condizione
Sono il Thread 1 ed ho aggiornato il contatore: 16
Sono il Thread 1 e sono bloccato sulla variabile condizione
Sono il Thread 2 ed ho aggiornato il contatore: 17
Sono il Thread 2 e sono bloccato sulla variabile condizione
Sono il Thread 1 ed ho aggiornato il contatore: 18
Sono il Thread 1 e sono bloccato sulla variabile condizione
Sono il Thread 2 ed ho aggiornato il contatore: 19
Sono il Thread 1 ed ho aggiornato il contatore: 20
Il valore finale del contatore e': 20
\end{lstlisting}
\captionof{lstlisting}{Output del Programma di cui sopra.}

\section{Esempio di utilizzo di una Variabile Atomica} \label{appendix:AtomicVarExample}
Si riporta un esempio di implementazione di un contatore tramite una Variabile di tipo Atomico, si noti l'assenza dei costrutti di Sincronizzazione sopra riportati.
\begin{lstlisting}[language=Java, numbers=left, stepnumber=1]
/*
 * File: AtomicVariable.java
 * Il Programma mira a mostrare le caratteristiche di
 * sincronizzazione delle Classi Atomiche in Java
 */

import java.util.concurrent.atomic.AtomicInteger;

public class AtomicVariable {
    // Inizializzo un contatore atomico
    private AtomicInteger counter = new AtomicInteger(0);
    
    // Main
    public static void main(String[] args) {
        new AtomicVariable();
    }

    // Costruttore dell'oggetto principale
    private AtomicVariable() {
        // Creo due thread che incrementino il contatore in maniera concorrente
        Thread t1 = new AtomicVariableThread(1, counter);
        Thread t2 = new AtomicVariableThread(2, counter);
        t1.start();
        t2.start();
        // Attendo la terminazione e stampo il risultato
        try {
            t1.join();
            t2.join();
        } catch(InterruptedException ie) {
            System.out.println("Thread interrotto inaspettatamente");
        }
        System.out.println("Il valore finale del contatore e': " + this.counter.get());
    }
}

class AtomicVariableThread extends Thread {
    // Identificatore
    private int threadNum;
    // Risorsa condivisa
    private AtomicInteger counter;
    
    // Costruttore
    protected AtomicVariableThread(int threadNum, AtomicInteger counter) {
        this.threadNum = threadNum;
        this.counter = counter;
    }
    
    // Implementazione del Metodo run()
    @Override
    public void run() {
        for(int i=0; i<10; i++) {
            System.out.println("Sono il Thread " + this.threadNum +
                " ed il valore del contatore e': " + this.counter.incrementAndGet());
        }
    }
}
\end{lstlisting}
\captionof{lstlisting}{Esempio di utilizzo di una variabile Atomica.}
\begin{lstlisting}[language=Bash]
$ java AtomicVariable
Sono il Thread 1 ed il valore del contatore e': 1
Sono il Thread 2 ed il valore del contatore e': 2
Sono il Thread 1 ed il valore del contatore e': 3
Sono il Thread 2 ed il valore del contatore e': 4
Sono il Thread 1 ed il valore del contatore e': 5
Sono il Thread 2 ed il valore del contatore e': 6
Sono il Thread 1 ed il valore del contatore e': 7
Sono il Thread 2 ed il valore del contatore e': 8
Sono il Thread 1 ed il valore del contatore e': 9
Sono il Thread 2 ed il valore del contatore e': 10
Sono il Thread 1 ed il valore del contatore e': 11
Sono il Thread 2 ed il valore del contatore e': 12
Sono il Thread 1 ed il valore del contatore e': 13
Sono il Thread 2 ed il valore del contatore e': 14
Sono il Thread 1 ed il valore del contatore e': 15
Sono il Thread 2 ed il valore del contatore e': 16
Sono il Thread 1 ed il valore del contatore e': 17
Sono il Thread 2 ed il valore del contatore e': 18
Sono il Thread 1 ed il valore del contatore e': 19
Sono il Thread 2 ed il valore del contatore e': 20
Il valore finale del contatore e': 20
\end{lstlisting}
\captionof{lstlisting}{Output del Programma di cui sopra.}